\section{Training Setup}

\subsection{Objective and Data Loading}

Training is formulated as next-token prediction. The xLSTM trainer loads MIDI paths from explicit split-list files, tokenizes and chunks them with MidiTok, wraps the chunks in \verb|DatasetMIDI|, and uses a \verb|DataCollator| with copied inputs as labels and shifted labels. Padding labels are ignored with index \verb|-100| in the cross-entropy loss.

The current block size is 1024 tokens. With tokenizer 11 and batch size 8, the training code reports 4,235 batches per epoch and 12,705 total training steps for three epochs.

\subsection{Optimization, Scheduling, and Logging}

The optimizer is AdamW using the xLSTM model's weight-decay parameter grouping. The base batch-8 training config uses:
\begin{itemize}
	\item seed 42;
	\item per-device train and validation batch size 8;
	\item learning rate / maximum learning rate \(3 \times 10^{-4}\);
	\item minimum learning rate \(3 \times 10^{-5}\);
	\item warmup for 200 steps;
	\item cosine decay until step 4000;
	\item 3 epochs;
	\item evaluation, checkpointing, and saving every 1000 steps;
	\item logging every 50 steps;
	\item W\&B logging enabled for batch-specific configs.
\end{itemize}

The trainer supports both legacy warmup/cosine fields and explicit multi-phase learning-rate schedules. The newer LR sweep configs use schedules such as linear warmup, cosine decay, and a final constant phase. Checkpoints are written to the run directory, and the best validation-loss checkpoint is also stored under \verb|best/checkpoint.pt|.

The xLSTM wrapper command is:
\begin{verbatim}
scripts/train_xlstm.sh <train_list> <val_list> <model_cfg> \
  <tokenizer_cfg> <output_dir> [train_cfg]
\end{verbatim}

